\documentclass[11pt,a4paper]{article}
\usepackage[utf8]{inputenc}
\usepackage[T1]{fontenc}
\usepackage[scaled]{helvet}
\renewcommand{\familydefault}{\sfdefault}
\usepackage[spanish,es-noshorthands]{babel}
\usepackage[a4paper,margin=20mm]{geometry}
\usepackage{xcolor}
\usepackage{tikz}
\usetikzlibrary{positioning,arrows.meta,calc}
\usepackage[most]{tcolorbox}
\usepackage[hidelinks]{hyperref}
\pagestyle{empty}
\setlength{\parindent}{0pt}

\definecolor{tinta}{HTML}{1D2330}
\definecolor{suave}{HTML}{5D6576}
\definecolor{acento}{HTML}{1F5F8B}
\definecolor{guinda}{HTML}{7A1F2B}
\definecolor{linea}{HTML}{DCD9D0}

% Enlace del video: reemplazar cuando esté subido a Drive.
\newcommand{\EnlaceVideo}{[pegar aquí el enlace del video en Drive]}

\begin{document}
\color{tinta}
\hyphenpenalty=10000\exhyphenpenalty=10000

%==================================================================== Hoja 1
{\small\color{suave} Auditoría de Sistemas · Trabajo de Investigación Formativa\\
Escuela Profesional de Ingeniería de Sistemas · Universidad Nacional de San Agustín de Arequipa}

{\color{guinda}\rule{\linewidth}{1.2pt}}

{\raggedright
\vspace*{38mm}
{\fontsize{30}{34}\selectfont\bfseries Agente de Terceros (DPAs)\par}
\vspace{4mm}
{\large\color{suave} Sistema de gestión de privacidad multiagente basado en ISO/IEC 27701:2025\par}

\vspace{16mm}
{\fontsize{24}{28}\selectfont\bfseries\color{guinda} Grupo 1\par}

\vspace{10mm}
{\large\bfseries\color{guinda} INTEGRANTES\par}
\vspace{3mm}
{\Large
\begin{tabular}{@{}l@{}}
Adrian Issac Mamani Quispe\\[3mm]
Dylan Antonio Zúñiga Huarca\\[3mm]
Jeferson Joao Basurco Casani
\end{tabular}\par}

\par}
\vfill
{\small\color{suave} Código fuente: \url{https://github.com/S-kkipie/agente-terceros-dpa}}

\newpage
%==================================================================== Hoja 2
{\small\color{suave} Agente de Terceros (DPAs) · Grupo 1}

{\color{guinda}\rule{\linewidth}{1.2pt}}

\vspace{2mm}
{\LARGE\bfseries Arquitectura\par}
\vspace{8mm}

\begin{center}
\begin{tikzpicture}[font=\small, node distance=6mm,
  caja/.style={draw=suave, rounded corners=3pt, align=center, minimum height=13mm, fill=white, line width=.7pt},
  ext/.style={caja, fill=black!4, draw=black!35},
  pri/.style={caja, fill=acento!10, draw=acento, line width=1.3pt},
  hum/.style={caja, fill=guinda!8, draw=guinda},
  fl/.style={-{Stealth[length=2.6mm]}, draw=suave, line width=.8pt},
  et/.style={font=\scriptsize\color{suave}, fill=white, inner sep=1.5pt}]

\node[ext, text width=46mm] (web) at (0,0) {\textbf{Interfaz web}\\ Panel del oficial de datos};
\node[ext, text width=46mm] (orq) at (5.6,0) {\textbf{Otros agentes del PIMS}\\ Incidentes, orquestador};
\node[hum, text width=46mm] (hitl) at (11.2,0) {\textbf{Aprobación humana}\\ Asesoría Legal / DPO};

\node[pri, text width=158mm, minimum height=16mm] (ag) at (5.6,-2.6)
  {\textbf{\large Agente de Terceros (DPAs)}\\[1mm]
   \ttfamily\footnotesize MCP: list\_processors · check\_dpa\_status · evaluate\_vendor · processors\_in\_breach};

\draw[fl] (web.south) -- node[et]{API REST} (web.south |- ag.north);
\draw[fl] (orq.south) -- node[et]{MCP} (orq.south |- ag.north);
\draw[fl] (hitl.south |- ag.north) -- node[et]{decisiones pendientes} (hitl.south);

\node[caja, text width=46mm] (rev) at (0,-5.2) {\textbf{Revisor de contratos}\\ 13 cláusulas · citas verificadas};
\node[caja, text width=46mm] (rie) at (5.6,-5.2) {\textbf{Evaluación de riesgo}\\ probabilidad × impacto};
\node[caja, text width=46mm] (bd) at (11.2,-5.2) {\textbf{Base de datos}\\ proveedores · contratos · bitácora};
\draw[fl] (rev.north |- ag.south) -- (rev.north);
\draw[fl] (rie.north |- ag.south) -- (rie.north);
\draw[fl] (bd.north |- ag.south) -- (bd.north);

\node[ext, text width=46mm] (faiss) at (0,-7.8) {\textbf{Normas indexadas}\\ FAISS + embeddings en Ollama (local)};
\node[ext, text width=104mm] (llm) at (8.4,-7.8) {\textbf{Modelo de lenguaje}\\ Gemini 2.5 Flash vía OpenRouter · respaldo automático (Gemini u Ollama)};
\draw[fl] (rev) -- (faiss);
\draw[fl] (rev.south east) -- (llm.north west);
\end{tikzpicture}
\end{center}

\vspace{6mm}
\small
\begin{tabular}{@{}p{36mm}p{126mm}@{}}
\textbf{Interfaz web} & Panel del oficial de datos: proveedores, revisión de contratos, aprobaciones, brechas, consulta normativa y bitácora.\\[1.5mm]
\textbf{Otros agentes} & Usan las herramientas del agente por MCP; por ejemplo, el Agente de Incidentes ante una brecha.\\[1.5mm]
\textbf{Revisor de contratos} & Contrasta cada contrato de encargo con 13 cláusulas mínimas; cada cita debe existir literal en el contrato.\\[1.5mm]
\textbf{Evaluación de riesgo} & Datos sensibles, volumen, transferencias, subencargados, incidentes y estado del contrato.\\[1.5mm]
\textbf{Normas indexadas} & Ley 29733, D.S. 016-2024-JUS, RGPD, Directrices 07/2020 del CEPD y cláusulas tipo (UE) 2021/915.\\[1.5mm]
\textbf{Aprobación humana} & Todo contrato nuevo y todo cambio de subencargado queda pendiente hasta que Legal o el DPO decide.
\end{tabular}

\vfill
\begin{tcolorbox}[colback=acento!5, colframe=acento, boxrule=1pt, arc=3pt, left=5mm, right=5mm, top=3mm, bottom=3mm]
{\large\bfseries\color{guinda} VIDEO DEL AGENTE FUNCIONANDO}\\[2mm]
{\large\bfseries\color{acento} \EnlaceVideo}
\end{tcolorbox}

\end{document}
